\section{Conclusions}
\label{sec:conclusions}

We have identified CsInH$_3$ in the cubic Pm$\bar{3}$m perovskite
structure as a phonon-mediated superconductor with
$T_c = 214$~K at $P = 3$~GPa ($\mu^* = 0.13$) after SSCHA anharmonic
corrections, matching the experimental $T_c = 203$~K of
H$_3$S~\cite{Drozdov2015} at a pressure reduced by a factor of~$30$.
The structure is thermodynamically viable
($E_{\mathrm{hull}} = 6$~meV/atom) and quantum-stabilized by hydrogen
zero-point motion at $3$~GPa, with all SSCHA phonon frequencies real.
However, the $\lambda$--$\omega_{\mathrm{log}}$ tradeoff intrinsic to
the MXH$_3$ perovskite family imposes a $T_c$ ceiling of
$\sim 214$--$234$~K within the Migdal-Eliashberg framework: anharmonic
corrections invariably reduce $\lambda$ by $30$--$36\%$ as phonons
soften toward the stability boundary, precluding room-temperature
superconductivity in this chemical class.

The $30\times$ pressure reduction --- from diamond-anvil-cell conditions
to the multi-anvil regime --- transforms the experimental accessibility
of high-$T_c$ hydride superconductivity, enabling bulk transport,
specific heat, and neutron scattering measurements on millimeter-scale
samples. Three validation steps are essential before experimental
pursuit: (i)~full DFPT~+~EPW calculations to resolve the
$\omega_{\mathrm{log}}$ discrepancy with
Du~\textit{et~al.}~\cite{Du2024}, which could
revise $T_c$ downward to $140$--$180$~K; (ii)~anisotropic Eliashberg
calculations to bound multi-gap effects on the perovskite Fermi surface;
and (iii)~experimental synthesis at $3$--$5$~GPa with \textit{in~situ}
transport measurements. More broadly, the perovskite $T_c$ ceiling
motivates exploration of alternative low-pressure hydride architectures
--- Laves-phase, Heusler, and modified clathrate families --- that may
circumvent the $\lambda$--$\omega_{\mathrm{log}}$ constraint and extend
near-ambient-pressure superconductivity beyond $200$~K.

%% CITATIONS NEEDED (for gpd-bibliographer):
%% - Drozdov2015 -- Drozdov et al., H3S 203 K, Nature 525, 73 (2015)
%% - Du2024 -- Du et al., Advanced Science 11, 2408370 (2024), MXH3 perovskites
